\Needspace{12\baselineskip}
\section{Current tests for smooth positivity}
\label{sec:smooth-cone-current-duality}

Let $X$ again be a connected compact complex manifold of dimension $n\geq1$. Testing against all positive $\ddc$-closed currents detects smooth positivity of a background. Hahn--Banach separation gives the exact threshold, with zero as the nef boundary. The full current set is needed here: the smooth tests in Lamari duality detect pseudoeffectivity.

Fix a Hermitian metric $h$ and use $\mathscr G_h$ from \cref{sec:gauduchon-linear-duality}. That set is compact for $n\geq2$. For $n=1$, elliptic regularity and the maximum principle give $\mathscr G_h=\{1/\int_Xh\}$, so compactness holds in this case as well.
\Needspace{6\baselineskip}
For any smooth real $(1,1)$-form $\theta$, define
\[
 m_h(\theta)=\min_{T\in\mathscr G_h}\int_X\theta\wedge T.
\]
The minimum exists by weak compactness, and it is finite since $-Ch\leq\theta\leq Ch$ for some $C$.

\begin{theorem}\label{thm:smooth-cone-current-duality}\index{Cone duality!smooth forms and currents}
The background $\theta$ has a smooth positive representative if and only if $m_h(\theta)>0$. Moreover $\theta$ is nef if and only if $m_h(\theta)\geq0$, and
\[
 m_h(\theta)=\sup\{a\in\mathbb R:\ \exists u\in C^\infty(X,\mathbb R),
                                  \ \theta+\ddc u>ah\}.
\]
\end{theorem}
\begin{proof}
If $\theta+\ddc u\geq ah$ for $a>0$, then distributional integration by parts gives $\int_X\theta\wedge T\geq a$ for every normalized test current. This proves one implication.

For the converse, work in the real Fr\'echet space of smooth $(1,1)$-forms. Let $\mathcal P$ be its open convex cone of strictly positive forms, and let $L=\ddc C^\infty(X,\mathbb R)$. If the affine space $\theta+L$ is disjoint from $\mathcal P$, the Hahn--Banach separation theorem for an open convex set and a disjoint affine space gives a nonzero continuous real linear functional $T$ with
\[
 T(\theta+\ell)\leq T(\eta)\qquad(\ell\in L,\ \eta\in\mathcal P).
\]
Varying $\ell$ by arbitrary real multiples shows $T(L)=0$. Scaling any $\eta\in\mathcal P$ to zero shows $T(\theta)\leq0$, while scaling it to infinity shows $T(\eta)\geq0$. Thus $T$ is a positive current: every smooth semipositive test form is a limit of strictly positive forms. Positivity also makes it a current of order zero. The equality $T(L)=0$ is precisely $\ddc T=0$. Since $T\ne0$ and $h>0$, its trace $T(h)$ is positive. Normalizing produces an element of $\mathscr G_h$ with nonpositive pairing against $\theta$. This contradicts $m_h(\theta)>0$ and proves the converse.

The identity $m_h(\theta-ah)=m_h(\theta)-a$ and the equivalence just proved show that the set inside the supremum is exactly $(-\infty,m_h(\theta))$. Finally, nefness means that $\theta+\varepsilon h$ has a smooth positive representative for every $\varepsilon>0$. This is equivalent to $m_h(\theta)+\varepsilon>0$ for all such $\varepsilon$, or $m_h(\theta)\geq0$.
\end{proof}

\Needspace{5\baselineskip}
For $n\geq2$, let $\mathscr C_h$ be the weak closure of the smooth positive normalized tests and let $\kappa_h(\theta)$ be their linear positivity threshold, as in \cref{sec:gauduchon-linear-duality}. The inclusion $\mathscr C_h\subseteq\mathscr G_h$ gives
$m_h(\theta)\leq\kappa_h(\theta)$.
The larger test set detects smooth positivity and nefness; the weak
closure of the smooth tests detects positivity and bigness of current
representatives. The inequality can be strict, as
\cref{ex:gauduchon-current-smoothing-obstruction} shows.

\begin{proposition}\label{prop:nef-annihilator-determinant}
Let $\theta$ be nef and suppose $\lambda_\mu(\theta)>0$. If $T\in\mathscr G_h$ satisfies $\int_X\theta\wedge T=0$, then its determinant-weighted measure vanishes.
\end{proposition}
\begin{proof}
In dimension one the hypotheses cannot hold, since
$\int_X\theta\wedge T=\lambda_\mu(\theta)/\int_Xh>0$.
Suppose $n\geq2$, and write $D_\mu(T)$ for the determinant functional of \cref{sec:weak-gauduchon-duality}, with $h$ as reference metric and volume form $\mu$. If $\eta\in\mathscr P$, its solution form satisfies $\Omega^n=\lambda_\mu(\eta)^n\mu$. Test the infimum in \cref{lem:weak-gaud-determinant-dual} with $\Omega/\lambda_\mu(\eta)$. Stokes gives
\[
 \int_X\eta\wedge T=\int_X\Omega\wedge T
              \geq\lambda_\mu(\eta)D_\mu(T).
\]
Apply this to $\eta=\theta+\varepsilon h$. The left side is $\varepsilon$, while \cref{thm:gaud-volume-nef-limit} gives $\lambda_\mu(\eta)\to\lambda_\mu(\theta)>0$. Thus $D_\mu(T)=0$. The determinant-weighted measure is positive and has this total mass, so it vanishes.
\end{proof}

A nef boundary with positive variational volume therefore has annihilating currents whose coefficient matrices are rank deficient almost everywhere. Their trace measures need not be supported on analytic subsets merely because these determinants vanish.

\Needspace{10\baselineskip}
\subsection{Maximal time of the Chern--Ricci flow}
\label{subsec:chern-ricci-maximal-time}

A smooth Chern--Ricci flow gives positive representatives along an affine ray. Conversely, a positive representative at a later time supplies a uniformly positive reference path and prevents an earlier singularity. The current criterion then computes the lifetime. Let $\omega_0$ be the initial Hermitian metric. For a Hermitian form $\omega$ on $X$, use the conventional Chern--Ricci form
\[
 \operatorname{Ric}(\omega)=-\mathrm{i}\partial\bar\partial\log\omega^n
                          =-2\pi\ddc\log\omega^n,
\]
interpreted in local coordinates; coordinate changes do not affect the form. This fixes the time normalization of the flow $\partial_t\omega=-\operatorname{Ric}(\omega)$. Put $\chi=\operatorname{Ric}(\omega_0)$ and $\theta_t=\omega_0-t\chi$.

\begin{theorem}\label{thm:chern-ricci-maximal-time}\index{Chern--Ricci flow!maximal time}
The Chern--Ricci flow starting from $\omega_0$ has a unique maximal smooth solution on $[0,\tau)$, where
\[
 \tau=\sup\{t\geq0:\ \theta_t\text{ has a smooth positive representative}\}.
\]
\end{theorem}
\begin{proof}
The admissible times form an interval containing zero: positive representatives at two times give positive representatives at the intervening times by convex combination. Positivity is also open in time. Fix an admissible $T>0$ and choose a smooth real function $a$ such that $\theta_T+\ddc a>0$. The reference forms
\[
 \gamma_t=\theta_t+\frac tT\ddc a
          =\left(1-\frac tT\right)\omega_0
             +\frac tT(\theta_T+\ddc a),\qquad 0\leq t\leq T,
\]
are uniformly positive, with all their derivatives uniformly bounded. Put
\[
 \Omega=\mathrm{e}^{a/(2\pi T)}\omega_0^n.
\]
Then $\dot\gamma_t=-\operatorname{Ric}(\Omega)$, with the Ricci form of a volume form defined by the same local expression. The scalar equation is
\begin{equation}\label{eq:chern-ricci-scalar-flow}
 \dot u=2\pi\log\frac{(\gamma_t+\ddc u)^n}{\Omega},
 \qquad \gamma_t+\ddc u>0,\qquad u(0)=0.
\end{equation}
Applying $\ddc$ shows that $\omega(t)=\gamma_t+\ddc u(t)$ solves the metric flow. Conversely, integrating $\dot u=2\pi\log(\omega(t)^n/\Omega)$ from zero recovers this identity from any metric solution. The linearization of the scalar equation is $\partial_t-2\pi\operatorname{tr}_{\omega(t)}\ddc$, a strictly parabolic operator. Classical short-time existence for smooth strictly parabolic equations therefore supplies a solution. Two scalar solutions have a difference satisfying a linear parabolic equation, whose positive principal coefficients are the averages of the inverse metrics along the segment between them. The maximum principle proves uniqueness.

Suppose this solution has maximal time $T_*<T$. We establish estimates independent of $t<T_*$ to continue it. For a large fixed $C$, the functions $Ct$ and $-Ct$ are respectively a supersolution and a subsolution of \eqref{eq:chern-ricci-scalar-flow}: at a spatial extremum of $u$ the complex Hessian has the required sign, and $\gamma_t^n/\Omega$ is bounded above and away from zero. The maximum principle gives $|u|\leq CT$.

Set $v=\dot u$ and $\mathcal L=2\pi\operatorname{tr}_{\omega(t)}\ddc$. Differentiation gives
\[
 (\partial_t-\mathcal L)v=2\pi\operatorname{tr}_{\omega(t)}\dot\gamma_t,
 \qquad
 (\partial_t-\mathcal L)u=v-2\pi n+2\pi\operatorname{tr}_{\omega(t)}\gamma_t.
\]
Since $\gamma_t$ is affine in time, these imply
\begin{align*}
 (\partial_t-\mathcal L)(t v-u-2\pi nt)
     &=-2\pi\operatorname{tr}_{\omega(t)}\omega_0<0,\\
 (\partial_t-\mathcal L)((T-t)v+u+2\pi nt)
     &=2\pi\operatorname{tr}_{\omega(t)}\gamma_T>0.
\end{align*}
The first expression is bounded above by its initial value, zero; the second is bounded below by $\min_X T v(0)$. The bound for $u$ and $T-t\geq T-T_*>0$ give a lower bound for $v$. The first inequality gives its upper bound on $[T_*/2,T_*)$, and the initial smooth solution bounds it on $[0,T_*/2]$. In particular, $\omega(t)^n/\Omega=\mathrm{e}^{v/(2\pi)}$ is bounded above and away from zero.

\Needspace{7\baselineskip}
We next bound the trace, retaining the time derivative of the reference metric. In the coordinates of \cref{lem:hermitian-trace-inequality}, write $g(t)$ for the matrix of $\gamma_t$ and $h=g+(u_{i\bar j})$ for that of $\omega(t)$. Put
\[
 S=\operatorname{tr}_g h,\qquad B=\operatorname{tr}_h g,
 \qquad F=\log\frac{\Omega}{\gamma_t^n}.
\]
Thus $\log(\det h/\det g)=v/(2\pi)+F$. The family $g(t)$ has uniformly bounded Chern torsion and curvature. At each fixed time \cref{lem:hermitian-trace-inequality} applies with these uniform constants. Direct differentiation of $S$ gives
\[
 \partial_t\log S
   =\frac{\Delta_g v+\operatorname{tr}_g\dot g}{S}
       -\frac{\operatorname{tr}(g^{-1}\dot g\,g^{-1}h)}{S}.
\]
The $\Delta_g v/S$ term cancels the corresponding term in $\mathcal L\log S$. If $\tau_k=\sum_iT_{ik}^i$ is the torsion trace of $g(t)$, the resulting inequality on $\{S\geq1\}$ is
\begin{equation}\label{eq:chern-ricci-parabolic-trace}
 (\partial_t-\mathcal L)\log S
 \leq-\frac{4\pi}{S^2}\operatorname{Re}(h^{k\bar q}\tau_k S_{\bar q})+CB+C.
\end{equation}
Indeed $\Delta_gF$ and $\operatorname{tr}_g\dot g$ are uniformly bounded, $S^{-1}\leq B/n^2$, and $-Cg\leq\dot g\leq Cg$ bounds the last displayed quotient. These bounds account for every added time-dependent term.

Choose a constant $C_0$ so that $1\leq s=u+C_0\leq D$ on the whole time interval, and maximize
\[
 Q=\log S-Au+s^{-1}
\]
on $X\times[0,t_0]$, where $t_0<T_*$ and $A\geq1$ will be fixed independently of $t_0$. A maximum at time zero gives the desired trace bound from the initial data. At a later maximum, $\partial\log S=(A+s^{-2})\partial u$. Since $|\tau|_h^2\leq CB$, Young's inequality bounds the absolute value of the first term on the right of \eqref{eq:chern-ricci-parabolic-trace} by
\[
 2\pi s^{-3}|\partial u|_h^2
       +\frac{C(A+1)^2s^3}{S^2}B.
\]
If $S^2\leq(A+1)^2D^3$, the trace is already bounded at the maximum. In the remaining case the last coefficient is at most $C$, and
\[
 \begin{split}
 0\leq(\partial_t-\mathcal L)Q
 &\leq CB+C-(A+s^{-2})(v-2\pi n+2\pi B)
                       -2\pi s^{-3}|\partial u|_h^2\\
 &\leq(C-2\pi A)B+C(A+1).
 \end{split}
\]
Choose $A$ so that $2\pi A>C+1$. This bounds $B$ at the maximum. The eigenvalues give $S\leq n(\det h/\det g)B^{n-1}$ for $n\geq2$, and $S=\det h/\det g$ for $n=1$. The determinant bound therefore bounds $S$ at the maximum in every case. The bound for $u$ transfers this estimate from $Q$ to $S$ everywhere, uniformly in $t_0$. The determinant lower bound then bounds the smallest eigenvalue away from zero. Thus $\omega(t)$ and $\gamma_t$ are uniformly equivalent for $t<T_*$.

Uniform ellipticity now gives the higher estimates by classical local parabolic regularity. In a coordinate ball, write the logarithmic determinant as a concave function of the complex part of the real Hessian. On the fixed compact set of positive matrices just obtained, its supporting affine functions have uniformly positive, bounded slopes. Their infimum extends the operator to a concave uniformly elliptic function of every real symmetric matrix. The spatial and time dependence enters through the smooth family $g(t)$ and the fixed density $\Omega$. The parabolic Evans--Krylov estimate gives uniform $C^{2+\alpha,1+\alpha/2}$ bounds away from the initial time. The original smooth solution covers a fixed initial interval. Differentiating the equation and applying parabolic Schauder estimates then gives uniform bounds for all spatial and time derivatives. This is the parabolic counterpart of the local extension used in \cref{cor:classical-higher-estimate}.

The bounds for the time derivatives make $u(t)$ Cauchy in every spatial $C^k$ norm as $t\uparrow T_*$. Its smooth limit still has a strictly positive solution form by uniform ellipticity. Short-time existence with this limit as initial datum extends the solution beyond $T_*$, a contradiction. Hence the flow exists up to every admissible time $T$, and uniqueness patches these solutions on $[0,\tau)$.

\Needspace{5\baselineskip}
Finally, integrating the metric flow gives
\[
 \omega(t)=\theta_t+\ddc\left(2\pi\int_0^t
                   \log\frac{\omega(s)^n}{\omega_0^n}\,\mathrm{d}s\right).
\]
Every time at which a smooth positive solution exists is therefore admissible. This proves maximality.
\end{proof}

Combining the continuation theorem with the current criterion gives a numerical description of the lifetime.

\Needspace{14\baselineskip}
\begin{corollary}\label{cor:chern-ricci-current-lifetime}
Let
\[
 b=\max_{T\in\mathscr G_{\omega_0}}\int_X\chi\wedge T.
\]
Then the maximal smooth existence time is
\[
 \tau=\begin{cases}1/b,&b>0,\\+\infty,&b\leq0.\end{cases}
\]
If $\tau<\infty$, the limiting background $\theta_\tau$ is nef and there is a normalized current $T_*$ with $\int_X\theta_\tau\wedge T_*=0$. If also $\lambda_\mu(\theta_\tau)>0$, every such annihilating current has zero determinant-weighted measure.
\end{corollary}
\Needspace{6\baselineskip}
\begin{proof}
For $t\geq0$, normalization of the tests gives the exact identity
\[
 m_{\omega_0}(\theta_t)=1-tb.
\]
The smooth-positivity theorem and \cref{thm:chern-ricci-maximal-time} prove the formula for $\tau$. When $b>0$, the maximum defining $b$ is attained by weak compactness. Its maximizer $T_*$ annihilates $\theta_\tau$. The same identity gives $m_{\omega_0}(\theta_\tau)=0$, hence nefness. The last assertion is \cref{prop:nef-annihilator-determinant}.
\end{proof}

\Needspace{8\baselineskip}
Along the flow, $\omega(t)=\theta_t+\ddc\varphi_t$ with a smooth potential. For every fixed $T\in\mathscr G_{\omega_0}$ one consequently has
\[
 \int_X\omega(t)\wedge T
       =1-t\int_X\chi\wedge T.
\]
These finite traces describe a boundary obstruction without any total Monge--Amp\`ere mass bound. They do not locate the singularities of the flow or prove smooth convergence away from an analytic set. Those geometric questions require local estimates in addition to the scalar pairings.

\Needspace{11\baselineskip}
\subsection{A positive-volume endpoint that is not big}
\label{subsec:actual-flow-positive-volume-nonbig}

Backgrounds reached at the maximal time of a Chern--Ricci flow are not better behaved, as far as the relation between positive volume and bigness is concerned. The radial background of \cref{sec:radial-positive-volume-nonbig} occurs, up to a smooth $\ddc$ change, as the limiting background of a Chern--Ricci flow at its maximal time. The normalized currents annihilating it can also be described explicitly.

We keep the notation of \cref{sec:radial-positive-volume-nonbig}: $X=\mathbb P^m\times\mathbb P^1$ with $n=m+1\geq3$, the normalized Fubini--Study forms $\omega$ and $w$ pulled back from the two factors, the nonnegative function $a=x/t$ on the first factor, which vanishes only at $p$, the background $\theta=\omega+aw$ and the probability form $\mu=\omega^m\wedge w$. Put $Z=\{p\}\times\mathbb P^1$, and start the flow from the Hermitian metric
\begin{equation}\label{eq:radial-flow-initial-metric}
 \omega_0=(n+1)\omega+(a+2)w.
\end{equation}

\begin{theorem}\label{thm:actual-flow-positive-volume-nonbig}\index{Chern--Ricci flow!nonbig positive-volume endpoint}
With the normalization of the Ricci form in \cref{subsec:chern-ricci-maximal-time}, the maximal smooth existence time of the Chern--Ricci flow starting from $\omega_0$ is $\tau=1/(2\pi)$. The limiting background $\alpha_\tau=\omega_0-\tau\operatorname{Ric}(\omega_0)$ is nef and not big, and
\[
 \alpha_\tau=\theta+\ddc\log(a+2),\qquad
 \mathcal V_\mu(\alpha_\tau)=n.
\]
Its variational volume is positive for every smooth strictly positive probability form on $X$.
\end{theorem}
\begin{proof}
The initial metric is strictly positive, and
\[
 \omega_0^n=n(n+1)^m(a+2)\mu.
\]
With our normalization, the Ricci forms of the normalized Fubini--Study metrics are $2\pi n\omega$ and $4\pi w$. Consequently
\[
 \operatorname{Ric}(\omega_0)
   =2\pi\bigl(n\omega+2w-\ddc\log(a+2)\bigr).
\]
Write $s$ for the flow time and put $r=2\pi s$. The affine background is
\begin{equation}\label{eq:radial-flow-affine-background}
 \alpha_s=(n+1-nr)\omega+(a+2-2r)w
                          +r\ddc\log(a+2).
\end{equation}
For $0\leq r<1$, adding $\ddc$ of the smooth function $-r\log(a+2)$ gives a strictly positive representative: the horizontal coefficient is at least one and the vertical coefficient is at least $2(1-r)>0$. At $r=1$ we obtain the identity $\alpha_\tau=\theta+\ddc\log(a+2)$. Nefness, bigness and the variational volume are unchanged by this smooth $\ddc$ change, by the conventions of \cref{chap:gauduchon-variational-volume} and \cref{prop:gaud-volume-elementary}, so \cref{thm:radial-positive-volume-nonbig} shows that the endpoint background is nef, not big, and of volume $n$.

For $r>1$, choose a nonzero smooth nonnegative function $f$ on the base supported in $\{a+2-2r<0\}$, a neighborhood of $p$. Its pullback $S=f\omega^m$ is a closed semipositive $(n-1,n-1)$-form. For every quasi-plurisubharmonic function $v$, Stokes' theorem for currents gives
\[
 \int_X(\alpha_s+\ddc v)\wedge S
       =\int_{\mathbb P^m}f(a+2-2r)\omega^m<0,
\]
so $\alpha_s$ is not even pseudoeffective for $r>1$. At $r=1$, a smooth positive representative would contradict nonbigness. Thus $r=1$ is the supremum of the admissible times, and \cref{thm:chern-ricci-maximal-time} gives the smooth flow on $[0,\tau)$ with exactly this lifetime; the reference representatives used above need not solve the flow equation. Finally, \cref{prop:gaud-volume-density-comparison} gives positivity of the volume for every smooth positive probability form.
\end{proof}

\begin{proposition}\label{prop:radial-flow-annihilator-support}\index{Gauduchon current!flow endpoint support}
The current $[Z]/2$ belongs to $\mathscr G_{\omega_0}$ and annihilates $\alpha_\tau$. Every $R\in\mathscr G_{\omega_0}$ with $\int_X\alpha_\tau\wedge R=0$ has its trace probability $\sigma_R=\omega_0\wedge R$ supported on $Z$. In the coefficient-matrix convention of \cref{sec:weak-gauduchon-duality}, its matrix is
\[
 P_R=n\Pi_{\mathrm{vert}}\qquad\sigma_R\text{-almost everywhere},
\]
where $\Pi_{\mathrm{vert}}$ is the $\omega_0$-orthogonal projection onto the tangent line of the fiber $\{p\}\times\mathbb P^1$. In particular, $P_R$ has rank one, and the determinant measure of $R$ vanishes for every smooth positive density.
\end{proposition}
\begin{proof}
Only smooth forms need to be restricted to $Z$, and their restrictions satisfy $\omega_0|_Z=2w$ and $\alpha_\tau|_Z=0$. This proves that $[Z]/2$ is normalized and annihilates $\alpha_\tau$.

Let $R$ be a general annihilator. Since $\alpha_\tau$ and $\theta$ differ by $\ddc$ of a smooth function, Stokes' theorem gives $\int_X\theta\wedge R=0$. The measure $\theta\wedge R$ is nonnegative, so it vanishes. On every compact set $K\subset X\setminus Z$, the explicit form $\theta$ satisfies $\theta\geq c_K\omega_0$ for some $c_K>0$. Thus $\sigma_R(K)=0$, and a countable exhaustion by compact sets shows that $\sigma_R$ is supported on $Z$.

Relative to $\omega_0$, the matrix of $R$ is characterized, for smooth test forms $\beta$ with matrix $B_\beta$, by
\[
 \beta\wedge R=\frac1n\operatorname{tr}(B_\beta P_R)\sigma_R,
 \qquad P_R\geq0,\qquad\operatorname{tr}P_R=n.
\]
At points of $Z$, the form $\theta$ is positive on the horizontal $m$-plane and vanishes on the vertical line. The vanishing of $\theta\wedge R$ therefore forces all horizontal diagonal entries of $P_R$ to vanish, and the inequality $|(P_R)_{ij}|^2\leq(P_R)_{ii}(P_R)_{jj}$ for positive matrices removes their rows and columns. The remaining vertical entry is $n$, which proves the formula and the vanishing of the determinant.
\end{proof}

For each such annihilator $R$, the trace identity at the end of \cref{subsec:chern-ricci-maximal-time} becomes
\[
 \int_X\omega(s)\wedge R=1-2\pi s,
 \qquad 0\leq s<1/(2\pi).
\]
The support statement uses the explicit semipositive representative $\theta$ of the endpoint class; the vanishing of the determinant alone would not give it.

